
\chapter{INTRODUCTION}\label{intro}
\vspace{-3em}

This chapter discusses ....



\section{Research  Background}
...


\section{Thesis Objectives}

The following are the objectives of this thesis: 
\begin{enumerate}
    \item To develop a model that will detect green space from aerial images.

    \item To propose a framework architecture for predicting the green space of a location in the next five years.

    \item To perform a comparative analysis of the outcomes of the trained models.
\end{enumerate}


\section{Methodological Overview}
To achieve the above objectives, this research was conducted in two phases - segmentation and prediction. Firstly, a literature study was conducted to understand how green space is being measured, the approaches followed to segment green spaces from images, what studies have been conducted so far related to forecasting green spaces in a designated area, and the methodology framework that has been defined. Secondly, the models and techniques followed for the segmentation and prediction of green space have been proposed. Finally, a comparative analysis was conducted to evaluate the performance of the conducted approaches and the best-performed system is proposed for predicting green spaces of a location from historical aerial images. \


\section{Scope of The Thesis}

The thesis primarily focused on developing a system that can predict the green space of a given location based on the historical aerial images of 15 years. The models were trained on aerial images. Therefore, the system is not designed for ground-level view. Moreover, images with a gap of 5 years were used for training the model. So, the system requires images with a gap of 5 years to give appropriate results. Additionally, the system is designed to give a numerical value (GVI difference value) as output.  \
 

\section{Organization of The Thesis}

The remainder of this research is organized as follows. 
The theoretical background and related works are briefly introduced in Chapter 2. 

The methodology which was followed throughout this research is presented in Chapter 3. The overall work has been divided into two phases (segmentation and prediction). Methodology of both the phases and the final proposed framework are described in brief.

 The overall evaluation of the system is discussed in Chapter 4. The evaluation consists of result analysis and comparative analysis of the functionalities and representation of both phases of the proposed framework. 
 
 The engineering considerations, challenges and remedies involved in this thesis have been discussed in Chapter 5.

 Finally, the discussion and conclusion are discussed in chapter 6. The discussion includes the thesis overview, contribution, research outcome, limitations, and future works.

